% Source: https://artofproblemsolving.com/community/c3h2446745_150th_msm_post
% Stable identity: aops-topic-2446745-q01
Let the radius be $r$; then the height is $h=2r$, and
\[V=\dfrac13\pi r^2h=\dfrac23\pi r^3=486\pi\quad\Longrightarrow\quad r^3=729,\quad r=9,\quad h=18.\]
The slant height is $\ell=\sqrt{r^2+h^2}=\sqrt{81+324}=\sqrt{405}=9\sqrt5$. The total surface area is the base plus the lateral surface:
\[\pi r^2+\pi r\ell=81\pi+81\sqrt5\,\pi=\boxed{81\pi\left(1+\sqrt5\right)}.\]

\begin{asy}
usepackage("amsmath");
size(5cm);
defaultpen(fontsize(10pt));
import graph;
real R=9,H=18,e=2.4;
path front=graph(new pair(real t){return (R*cos(t),e*sin(t));},pi,2pi);
path back=graph(new pair(real t){return (R*cos(t),e*sin(t));},0,pi);
draw(front,linewidth(0.8)); draw(back,dashed);
pair V=(0,H),O=(0,0),Q=(R,0);
draw((-R,0)--V--(R,0),linewidth(0.8));
draw(O--V,dashed); draw(O--Q,dashed);
draw((0,0.9)--(0.9,0.9)--(0.9,0));
dot(O); dot(V);
label("$h=18$",(0,6),E,fontsize(8pt)); label("$r=9$",(R/2,0),N,fontsize(8pt)); label("$\ell$",(R/2,H/2),NE);
label("$T$",V,N); label("$O$",O,W);
\end{asy}

(The lateral surface alone has area $81\sqrt5\,\pi$.)
